\documentclass{article}
\usepackage[utf8]{inputenc}
\usepackage[spanish]{babel}
\usepackage{amsmath}
\usepackage{geometry}
\geometry{a4paper, margin=1in}

\title{Resumen: Investigación de Operaciones}
\author{
    Alison Renata García Romero \\
    Uriel Garduño Salinas \\
    Irving Johan Gudiño Cruz \\
    Ángel Jafet Tenorio Zamora
}
\date{30 de septiembre de 2026}

\begin{document}

\maketitle

\section{Introducción}
La investigación de operaciones es una disciplina que permite traducir situaciones de la vida real (como la asignación de recursos escasos: horas de máquina, materias primas, presupuesto o personal) al lenguaje de las matemáticas para resolver problemas de manera óptima y traducirlos de regreso a decisiones de negocio[cite: 2]. El proceso consta de tres etapas fundamentales: entender la situación, definir variables, construir el modelo y, por último, interpretar la solución[cite: 2].

\section{Del problema a un modelo}
Un modelo de programación lineal (PL) responde tres preguntas: ¿qué se puede decidir?, ¿qué se quiere lograr? y ¿qué límites no se pueden cruzar?. Donde se siguen los siguientes puntos, como una receta general:
\begin{enumerate}
    \item Describir la decisión y las unidades[cite: 2].
    \item Definir una variable por cada decisión desconocida[cite: 2].
    \item Escribir la función objetivo (maximizar beneficio o minimizar costo)[cite: 2].
    \item Convertir cada límite de recursos, demanda o producción en una restricción[cite: 2].
    \item Añadir las condiciones de no negatividad[cite: 2].
    \item Revisar dimensiones, unidades y el significado de cada resultado (el paso que más se omite)[cite: 2].
\end{enumerate}

Forma general: optimizar $Z = c_1 x_1 + \dots + c_n x_n$, sujeto a restricciones lineales y $x_i \ge 0$. Lineal significa que cada variable aparece solo multiplicada por una constante, sin productos entre variables ni exponentes distintos de 1. Esto es lo que permite resolver el problema de forma sistemática y garantizada[cite: 2].

\section{Matrices para organizar información}
\subsection{Dimensiones y operaciones}
\begin{itemize}
    \item \textbf{Tamaño:} $A = (a_{ij})$ de $m \times n$ tiene $m$ filas y $n$ columnas; anótelo siempre antes de operar[cite: 2].
    \item \textbf{Suma y resta:} entrada por entrada, solo con matrices del mismo tamaño[cite: 2].
    \item \textbf{Producto AB:} requiere que columnas de $A =$ filas de $B$. Si $A$ es $m \times n$ y $B$ es $n \times p$, $AB$ es $m \times p$[cite: 2].
    \item \textbf{Transpuesta:} intercambia filas y columnas ($A^T$ es $n \times m$)[cite: 2].
    \item \textbf{Potencia $A^2$:} solo tiene sentido si $A$ es cuadrada[cite: 2].
\end{itemize}

\textbf{Construir una matriz a partir de una regla}, es fijar el tamaño; sustituir $i = 1$ y recorrer $j = 1\dots n$ para la primera fila; repetir con $i = 2\dots m$; comprobar posiciones. Si la regla es por casos, verificar primero a qué caso pertenece cada pareja $(i, j)$. En fórmulas con fracciones, calcular primero el numerador y luego dividir[cite: 2].

\textbf{Matriz de coeficientes de un sistema}, donde un sistema de $m$ ecuaciones se resume como $AX = B$, donde $A$ contiene solo los coeficientes (incluyendo ceros de variables ausentes), $X$ el vector de incógnitas y $B$ los términos independientes. El signo igual no forma parte de $A$. Se puede resolver con la matriz aumentada $[A \mid B]$ y eliminación de Gauss, usando tres operaciones: intercambiar filas, multiplicar una fila por un escalar no nulo y sumar a una fila un múltiplo de otra[cite: 2].

\section{Algoritmo Simplex (Maximización y Minimización)}
El algoritmo opera mediante tablas iterativas basándose en las siguientes reglas:
\begin{itemize}
    \item \textbf{Variable entrante:} En maximización entra el mayor costo reducido positivo ($C_j - Z_j > 0$); en minimización entra el más negativo ($C_j - Z_j < 0$)[cite: 2].
    \item \textbf{Variable saliente (Razón $\theta$):} Se calcula $\theta = b_i / a_{ik}$ únicamente para coeficientes positivos de la columna entrante. La menor razón determina la fila que sale[cite: 2].
    \item \textbf{Pivoteo:} Se normaliza la fila pivote dividiéndola entre el elemento pivote y se hacen ceros en el resto de la columna mediante operaciones elementales entre filas[cite: 2].
    \item \textbf{Criterio de parada:} El proceso concluye cuando ya no existen costos reducidos que mejoren la función objetivo[cite: 2].
\end{itemize}

Donde las \textbf{situaciones especiales} y la comprobación final se buscan que sean óptimos múltiples, es decir, una variable no básica con costo reducido cero permite múltiples combinaciones óptimas. Así como degeneración, coincidencia en las razones mínimas que puede generar bases con términos independientes nulos[cite: 2].

Pero también se encontrará \textbf{no factibilidad}, presencia de variables artificiales positivas al finalizar el algoritmo, indicando que el modelo no tiene solución factible. \textbf{No acotación}, es cuando la columna entrante no presenta coeficientes positivos, permitiendo un crecimiento indefinido[cite: 2].

\end{document}